\scanpage{6}\mspage{4}
car si
\[
a'\overset{u}{\unsure{\text{---}|\longrightarrow}}a\qquad b\xrightarrow{v}b'
\]
on a le diagramme
\diagram{\begin{tikzcd}[ampersand replacement=\&,column sep=large,row sep=large]
\mathrm{Ch}_{\infty}(a,b) \arrow[r,no head,"{\substack{\text{comp. avec}\\ v}}"] \arrow[d,"{\text{comp. avec}}"',"{u}"] \& \mathrm{Ch}_{\infty}(a,b') \arrow[d,"{\text{comp. avec }u}"]\\
\mathrm{Ch}_{\infty}(a',b) \arrow[r,no head,"{\substack{\text{comp. avec}\\v}}"] \& \mathrm{Ch}_{\infty}(a',b')
\end{tikzcd}}
qui est commutatif grâce à l'associativité pour
\[
\mathrm{Ch}_{\infty}(a',a),\ \mathrm{Ch}_{\infty}(a,b),\ \mathrm{Ch}(b,b').
\]

\uline{Corollaire} de cette \unsure{dépendance|disposition} bifonctorielle : Si $u:\deleted{\unsure{b}}a'\to\deleted{\unsure{b'}}a$, \inserted{$v:b\to b'$} \unsure{st|est} \unsure{isom.} dans $M_{\pi}$, alors la flèche correspondante
\[
\mathrm{Ch}_{\infty}(a,b)\longrightarrow\mathrm{Ch}_{\infty}(a',b')\qquad\text{dans Hot}
\]
est un \uline{iso}.

Les $\unsure{\overline{\mathrm m}\mathrm s|\overline{\mathrm m}\mathrm e\mathrm s}$ arguments devraient marcher si au lieu \unsure{des|du} $\mathrm{Hom}$ ext. à valeurs dans $\mathrm{Hot}$, on en prenait dans $\mathrm{Ind}\,\mathrm{Hot}$, ou dans $\mathrm{Cat}\,\mathrm{hot}$, ou dans $\mathrm{Ind}\,\mathrm{Cat}\,\mathrm{hot}$ -- je pense bien sûr aux $\mathrm{CH}_{\infty}(a,b)$ (dans la théorie de la localisation) dans $\mathrm{Ind}\,\mathrm{Cat}\,\mathrm{hot}$, qui \deleted{définissent} peuvent être regardés comme provenant d'un bifoncteur (où $M_{\pi}=M\Sigma^{-1}$)
\[
M_{\pi}^{\unsure{o}}\times M_{\pi}\ \text{---}\ \mathrm{Ind}\,\mathrm{Cat}\,\mathrm{hot}
\]
% archive sheet number: 5

